\documentclass[10pt,a4paper]{article}
\usepackage[spanish,es-nodecimaldot]{babel}
\usepackage[utf8]{inputenc}
\usepackage[T1]{fontenc}
\usepackage{lmodern}
\usepackage[top=1.1cm,bottom=1.1cm,left=1.5cm,right=1.5cm]{geometry}
\usepackage{amsmath,amssymb}
\usepackage[table]{xcolor}
\usepackage{enumitem}

\setlength{\parindent}{0pt}
\setlength{\parskip}{2pt}
\pagestyle{empty}

\AtBeginDocument{%
  \setlength{\abovedisplayskip}{4pt}%
  \setlength{\belowdisplayskip}{4pt}%
  \setlength{\abovedisplayshortskip}{2pt}%
  \setlength{\belowdisplayshortskip}{2pt}%
}

\begin{document}

{\large\bfseries Regresion curvilinea}
\vspace{2pt}

Ahora que pasa si al graficar el diagrama de dispersion vemos que los datos no siguen una trayectoria recta? queremos ajustar con linea recta pero el modelo no va a ser el mas adecuado ya q no representa la naturaleza del fenomeno. Es por esto q usamos un modelo curvilineo. Para lograr este modelo y que se comporte de manera lineal debemos aplicar una transformacion matematica a los datos originales y asi poder aplicar el metodo de los minimos cuadrados que ya veniamos trabajando.

\textbf{Exponencial}
\[
y = \alpha\cdot\beta^x
\;\Rightarrow\;
\log_{10}(y) = \log_{10}(\alpha\cdot\beta^x)
\;\Rightarrow\;
\log_{10}(y) = \log_{10}(\alpha) + x\cdot\log_{10}(\beta)
\]
\[
y^{*} = \log_{10}(y), \qquad a = \log_{10}(\alpha), \qquad b = \log_{10}(\beta)
\]
\[
\Rightarrow \quad y^{*} = a + b\cdot x
\qquad\qquad\qquad
\alpha = 10^a \qquad \beta = 10^b
\]

\textbf{Potencial}
\[
y = \alpha\cdot x^{\beta}
\;\Rightarrow\;
\log_{10}(y) = \log_{10}(\alpha) + \beta\cdot\log_{10}(x)
\]
\[
\Rightarrow \quad y^{*} = \log_{10}(y), \quad x^{*} = \log_{10}(x), \quad a = \log_{10}(\alpha), \quad B = \beta
\]
\[
\Rightarrow \; y^{*} = a + Bx^{*}
\qquad\qquad\qquad
\alpha = 10^a \qquad \beta = 10^b
\]

\textbf{Reciprocos}
\[
y = \frac{1}{\alpha+\beta\cdot x}
\;\Rightarrow\;
\frac{1}{y} = \alpha + \beta\cdot x
\]
\[
\Rightarrow \; y^{*} = \frac{1}{y} \qquad a=\alpha \quad b=\beta
\]
\[
\Rightarrow \; y^{*} = \alpha + \beta x
\]

\end{document}
